\documentclass[12pt,a4paper]{article}

% ============================================================
% Living manuscript source — update section by section.
% Compile with: pdflatex main.tex (run twice for references).
% ============================================================
\usepackage[margin=1in]{geometry}
\usepackage[T1]{fontenc}
\usepackage{lmodern}
\usepackage{microtype}
\usepackage{setspace}
\usepackage{graphicx}
\usepackage{booktabs}
\usepackage{longtable}
\usepackage{array}
\usepackage{amsmath,amssymb}
\usepackage{siunitx}
\usepackage{enumitem}
\usepackage{hyperref}
\usepackage{xcolor}
\usepackage{natbib}
\usepackage{caption}
\usepackage{float}
\hypersetup{colorlinks=true,linkcolor=blue,citecolor=blue,urlcolor=blue}
\setstretch{1.35}
\setlength{\parskip}{0.45em}
\setlength{\parindent}{0pt}

\title{\textbf{Integrated Urban Stormwater Quantity and Quality Modelling Using EPA-SWMM:}\\
A Critical Review of Methods, Evidence, Uncertainty and Sustainable Management}
\author{Adarsh Chourasiya\\
\small Affiliation and co-author details to be confirmed}
\date{Working manuscript v0.4 \quad | \quad 3 November 2026}

\begin{document}
\maketitle

\begin{center}
\textbf{Manuscript status: Living working draft --- not submission-ready}
\end{center}

\begin{quote}
\textbf{Research-integrity note.} This document is being developed section by section as the literature review and evidence appraisal progress. The current evidence register contains seven audited records, of which three have been appraised using accessible full texts and four remain at abstract or publisher-excerpt level. The candidate bibliography is not a completed systematic-review sample. No field observations, calibrated results, measured campus characteristics, or model outputs are claimed. The proposed MANIT Bhopal application is a framework, not a completed case study.
\end{quote}

\begin{abstract}
Urban stormwater runoff is both a hydrologic response to rainfall and a transport pathway for pollutants accumulated on urban surfaces. Integrated assessment therefore requires a coherent representation of rainfall--runoff generation, drainage hydraulics, pollutant mobilisation, water-quality routing, monitoring, model evaluation and management interventions. The U.S. Environmental Protection Agency Storm Water Management Model (EPA-SWMM) provides a widely used platform for single-event and continuous simulation of urban runoff quantity and quality. However, prediction credibility depends on catchment discretisation, rainfall and drainage data, pollutant-process assumptions, calibration observations, validation design and uncertainty treatment. This critical review examines the methodological links between runoff quantity and quality simulation, with emphasis on hydrology and hydraulics, pollutant buildup and washoff, event mean concentration (EMC), first-flush assessment, GIS representation, calibration and validation, uncertainty, and green infrastructure/low-impact development (GI/LID). Early evidence indicates that formal evaluation of water-quality predictions has been less common than runoff evaluation in the SWMM literature. Recent work also demonstrates that water-quality monitoring is constrained by the limited range of parameters measurable continuously and by site-specific sensor calibration and field-operation requirements. A recent large-catchment SWMM application illustrates the practical value and limitations of EMC-based nutrient modelling when pollutant observations are sparse. The review develops a decision-oriented framework for future application to a campus catchment in Bhopal, India. Such an application will require verified catchment and drainage data, local rainfall, flow observations and event-based water-quality measurements. The evidence synthesis remains in progress; conclusions are therefore provisional.
\end{abstract}

\textbf{Keywords:} urban stormwater; EPA-SWMM; runoff modelling; stormwater quality; buildup; washoff; event mean concentration; first flush; calibration; uncertainty; GIS; low-impact development.

\tableofcontents
\newpage

\section{Introduction}
Urbanisation changes the partitioning and routing of rainfall through the replacement or fragmentation of pervious surfaces by roofs, roads, parking areas and other impervious surfaces. These changes can increase runoff volume and peak discharge, shorten catchment response time and alter connectivity between pollutant source areas and receiving waters. During dry periods, suspended solids, organic matter, nutrients, metals and other contaminants may accumulate on urban surfaces. Rainfall can mobilise part of this material and transport it through overland flow and drainage networks. Stormwater management is consequently a coupled quantity-and-quality problem.

Urban stormwater models differ in their spatial and temporal resolution and in their representation of runoff generation, drainage hydraulics and water quality. EPA-SWMM is a dynamic rainfall--runoff model used for single-event and long-term simulation of runoff quantity and quality. Its subcatchment runoff component can generate runoff and pollutant loads, while its routing component transports flows through pipes, channels, storage and treatment elements \citep{rossman2022}. The model also includes water-quality buildup and washoff representations and low-impact-development controls. The availability of these components does not, by itself, guarantee accurate prediction: model performance depends on the quality of the input data, suitability of the process representation and adequacy of calibration and validation.

A key distinction in the literature is between hydrologic performance and water-quality performance. A reasonable runoff hydrograph fit does not demonstrate that pollutant concentrations, event loads or pollutographs are accurate. Niazi et al. \citeyearpar{niazi2017} reviewed SWMM model evaluation and found that water-quality endpoints were formally evaluated less often than runoff endpoints in the set of studies they examined. Their review also identified gaps in diffuse pollutant representation, contaminant fate and transport, parameter estimation, uncertainty analysis and GI/LID effectiveness. These findings support the need to evaluate both quantity and quality explicitly rather than treating water quality as an optional post-processing step.

Pollutant processes introduce additional complexity. Buildup formulations represent the accumulation of pollutant mass during dry periods, while washoff formulations represent mobilisation and removal during rainfall and runoff. The appropriate functional form and parameter values may vary by pollutant, surface type, rainfall event and catchment scale. For example, Gaut et al. \citeyearpar{gaut2019} reported different washoff-model performance for particulate and dissolved pollutants. Bonhomme and Petrucci \citeyearpar{bonhomme2017} cautioned that parameters calibrated at catchment scale may lose a straightforward physical interpretation, even when the underlying buildup/washoff equations are physically motivated at smaller scales.

Monitoring is therefore part of model design, not merely a final validation step. Stormwater quality varies between sites, between events and within individual events. A review of continuous monitoring methods found that only a limited set of parameters had been measured directly at high temporal resolution for extended periods, while many target pollutants were inferred from proxy sensor measurements \citep{razguliaev2024}. Flow observations are also necessary: pollutant load is calculated from concentration and discharge over time, not from concentration alone.

\subsection{Aim and objectives}
The aim of this paper is to critically review methods and evidence for integrated urban stormwater quantity and quality modelling using EPA-SWMM and to formulate a transparent framework for a future campus-catchment application in Bhopal, India.

The specific objectives are:
\begin{enumerate}[leftmargin=1.5em]
  \item to review the representation of rainfall--runoff processes and drainage hydraulics in SWMM;
  \item to examine pollutant sources, buildup, washoff, EMC and first-flush methods;
  \item to evaluate calibration, validation, sensitivity and uncertainty practices for quantity and quality endpoints;
  \item to assess the influence of GIS-derived catchment characteristics and drainage connectivity;
  \item to synthesise evidence on GI/LID interventions and their quantity--quality outcomes; and
  \item to identify the data and methodological requirements for a future monsoon-influenced campus application.
\end{enumerate}

\subsection{Review questions}
The review is guided by five connected questions:
\begin{enumerate}[leftmargin=1.5em]
  \item How are runoff generation and drainage hydraulics represented and evaluated in SWMM applications?
  \item How are pollutant buildup, washoff, EMC and first-flush behaviour represented, and how are these processes validated?
  \item Which monitoring, calibration, validation and uncertainty practices support credible simultaneous prediction of runoff and pollutant loads?
  \item How do spatial discretisation, land use, impervious connectivity and drainage topology influence model results?
  \item What evidence supports GI/LID scenario analysis, and what limits transfer to monsoon-influenced urban catchments?
\end{enumerate}

\section{Review methodology}
\subsection{Review design and current status}
The review is being conducted as a critical integrative review. It is not currently described as a completed systematic review. A preliminary candidate bibliography has been assembled, and the first evidence audit has verified bibliographic details for seven records. Three papers have been appraised from accessible full texts: Niazi et al. \citeyearpar{niazi2017}, Razguliaev et al. \citeyearpar{razguliaev2024}, and Bhusal and Muthukumaran \citeyearpar{bhusal2026}. Four additional records in the current audit remain at abstract or publisher-excerpt level. This distinction is retained because bibliographic verification, abstract assessment and full-text appraisal are different levels of evidence.

The next review stage will expand full-text retrieval and screening across SWMM hydrology and hydraulics, pollutant-process modelling, calibration and uncertainty, GIS representation, and GI/LID. The final search dates, database-specific queries, deduplication procedure, screening counts and exclusion reasons will be recorded before any systematic-review reporting is considered. PRISMA 2020 may be used as a reporting guideline if the completed workflow meets the requirements of a systematic review; it is not a substitute for conducting the search and screening process \citep{page2021}.

\subsection{Information sources and search concepts}
The intended information sources include Scopus, Web of Science, Engineering Village, ASCE Library, ScienceDirect, SpringerLink, Wiley, Taylor \& Francis, Crossref and supplementary citation searching. Search concepts will combine terms for (i) SWMM and urban drainage, (ii) runoff quantity and hydraulic routing, (iii) stormwater quality and pollutant loads, (iv) buildup and washoff, (v) EMC and first flush, (vi) calibration, validation and uncertainty, (vii) GIS and catchment discretisation, and (viii) GI/LID and sustainable drainage.

An example search string is:
\begin{quote}
(\texttt{SWMM OR "Storm Water Management Model"}) AND
(\texttt{stormwater OR "urban runoff"}) AND
(\texttt{quality OR pollutant OR buildup OR washoff OR EMC OR "first flush" OR calibration OR uncertainty OR LID})
\end{quote}
The exact search string will be adapted to the syntax of each database and archived with the date of execution.

\subsection{Eligibility and evidence handling}
Priority will be given to peer-reviewed journal articles and substantive review papers that provide evidence relevant to urban runoff quantity, drainage hydraulics, stormwater quality, pollutant processes, monitoring, model evaluation or GI/LID. Foundational studies will be retained where they clarify model structure or process representation. Unverifiable records, duplicate records and sources without a defensible connection to the review questions will be excluded or clearly labelled as background material.

For each eligible study, the extraction matrix will record the citation and DOI, study location, catchment area and land use, rainfall and monitoring design, model and version, pollutants, process representation, parameter sources, calibration and validation design, performance metrics, uncertainty analysis, GI/LID practices, limitations and transferability. A descriptive appraisal rubric will assess data adequacy, process transparency, independent validation, uncertainty treatment and reproducibility. No single composite quality score will be used unless its weighting is justified.

\section{Integrated stormwater quantity and quality processes}
\subsection{Rainfall--runoff transformation}
Rainfall excess is converted into surface runoff through interacting processes that include interception, infiltration, depression storage and surface routing. In SWMM, subcatchments are represented through effective parameters such as area, width, slope, imperviousness, surface roughness and depression storage. These parameters influence runoff volume, peak discharge and hydrograph timing. Their interpretation requires care: some are derived from GIS or field surveys, while others are effective parameters adjusted during calibration. A calibrated effective parameter should not automatically be described as a directly measured physical property.

The spatial distribution and connectivity of impervious surfaces are particularly important. Two catchments with the same total impervious fraction can produce different hydrographs if roofs, roads and paved areas connect differently to the drainage network. Subcatchment delineation and discretisation should therefore be guided by topography, land use, drainage connectivity and the intended management question. Excessive subdivision can increase parameter burden without improving identifiability; overly coarse aggregation can conceal important source-area and routing differences.

\subsection{Drainage hydraulics}
The drainage network routes runoff through junctions, conduits, channels, storage elements, pumps, regulators and outfalls. Network topology, conduit dimensions, invert elevations, roughness, inlet representation and boundary conditions affect flow depth, travel time, surcharge and flooding. The dynamic-wave routing option can represent backwater and surcharge effects, but its use requires reliable network geometry and appropriate computational settings. A hydraulically detailed model built from uncertain geometry may create an appearance of precision without corresponding predictive reliability.

For integrated quantity--quality modelling, hydraulic routing also determines the timing and dilution of pollutant concentrations. Errors in flow volume or travel time can propagate into pollutant-load and pollutograph errors. Hydrologic and hydraulic checks should therefore precede interpretation of water-quality performance.

\subsection{Pollutant sources and loads}
Urban stormwater pollutant sources vary by land use, surface material and activity. Road and parking runoff may carry suspended solids and particle-associated metals; landscaped areas, roofs and other surfaces may contribute different nutrient and organic signatures. Source inventories should be locally justified rather than copied wholesale from another city.

For concentration $C(t)$ and discharge $Q(t)$, the instantaneous pollutant mass flux is $C(t)Q(t)$. Event pollutant load is:
\begin{equation}
  L_{\mathrm{event}} = \int_{t_0}^{t_1} C(t)Q(t)\,dt,
  \label{eq:eventload}
\end{equation}
where $t_0$ and $t_1$ define the event interval. Equation~\ref{eq:eventload} makes clear why concentration alone is not sufficient to quantify exported mass. If concentration is reported in \si{mg/L} and discharge in \si{L/s}, the integrated result must be converted consistently to a mass unit.

\section{Pollutant buildup and washoff}
Buildup functions represent pollutant accumulation during dry-weather periods. Common conceptual forms include power, exponential and saturation-type functions. Washoff functions represent the mobilisation and removal of accumulated mass under rainfall and runoff, often using an exponential or rating-type relationship. These functions are parsimonious and can be useful where detailed mechanistic source data are unavailable, but their parameters require constituent-specific and site-specific justification.

Soonthornnonda et al. \citeyearpar{soonthornnonda2008} developed a washoff model with pollutant-specific coefficients, illustrating that solids, metals, nutrients and microbial indicators need not share a common response. Gaut et al. \citeyearpar{gaut2019} examined six pollutants across 126 events and 12 sites in Singapore; the abstract-level evidence in the current audit indicates stronger performance for particulate TSS than for dissolved nutrients. The full-text appraisal of that study remains pending, so its detailed metrics and validation design will not be reported here as established findings.

A scale problem also deserves attention. Bonhomme and Petrucci \citeyearpar{bonhomme2017} used a Bayesian calibration approach for a peri-urban catchment and argued that buildup/washoff parameters at catchment scale may not retain the physical meaning assigned to them from small-scale experiments. This raises a central issue for model interpretation: a good statistical fit does not necessarily prove that the calibrated parameters represent transferable physical processes.

For each pollutant, an SWMM application should document the constituent name and units, source-area assignment, buildup function, washoff function, parameter source, antecedent dry-period treatment, residual-mass assumption, simulation time step and calibration evidence. Where observations are sparse, parameters should be treated as uncertain and scenario results should be labelled accordingly.

\section{Event mean concentration and first-flush assessment}
The event mean concentration is a flow-weighted average concentration over a defined runoff event:
\begin{equation}
  EMC = \frac{\int_{t_0}^{t_1} C(t)Q(t)\,dt}
  {\int_{t_0}^{t_1} Q(t)\,dt}.
  \label{eq:emc}
\end{equation}
EMC is useful for event-load estimation and comparison between events, but it can conceal short-lived concentration peaks and within-event dynamics. A model that reproduces an event EMC may still fail to reproduce the timing and magnitude of the pollutograph.

First-flush analysis examines the relationship between cumulative pollutant mass fraction and cumulative runoff-volume fraction. It is sensitive to event boundaries, sampling interval, pollutant selection, catchment characteristics and the chosen first-flush criterion. A fixed criterion such as the first 20\% of runoff volume should not be assumed to apply universally. Monitoring should capture the rising limb and subsequent event periods, with flow and sample timestamps synchronised. Where resources permit, flow-weighted composite samples can be supplemented by discrete high-frequency samples to identify pollutograph dynamics.

\section{Monitoring, calibration, validation and uncertainty}
\subsection{Monitoring design}
Continuous monitoring can improve understanding of within-event stormwater quality dynamics, but the range of parameters measured directly at high temporal resolution remains limited. Razguliaev et al. \citeyearpar{razguliaev2024} reviewed 84 articles representing 56 independent studies across 16 countries and reported that only eight water-quality parameters had been measured directly on-site for extended periods in the reviewed literature. Many target pollutants were inferred from proxy variables, which requires site-specific calibration against laboratory measurements. Field operation, cleaning, sensor drift, power, data transmission and validation also affect data quality.

For a quantity--quality model, the minimum coherent monitoring package should pair rainfall, water level or discharge, and event-based water-quality samples. Rainfall and flow time steps should be sufficiently fine to capture event dynamics. Laboratory QA/QC, detection limits, sample preservation, missing-data handling and sensor calibration should be documented. Pollutant loads should be calculated only where concentration and flow observations can be aligned adequately in time.

\subsection{Calibration and validation}
Calibration estimates parameter values using observations, while validation evaluates the model on data not used in calibration. A defensible integrated evaluation should distinguish at least:
\begin{itemize}[leftmargin=1.5em]
  \item runoff volume, peak discharge and peak timing;
  \item hydrograph shape and event runoff response;
  \item pollutant event load and EMC;
  \item pollutograph timing, concentration peaks and event-to-event variability; and
  \item mass balance and, where relevant, receiving-water or treatment performance.
\end{itemize}
Hydrology-only calibration is insufficient when water-quality predictions are a primary research output. Calibration and validation should use multiple events where possible, and the validation set should be independent in time or event selection. Metrics such as NSE, KGE, $R^2$, RMSE, MAE and percent bias may be reported where appropriate, but no single statistic should be treated as proof of model adequacy. Graphical hydrograph and pollutograph comparisons, residual analysis and mass-balance checks are also necessary.

\subsection{Uncertainty and parameter transferability}
Uncertainty arises from rainfall inputs, catchment delineation, drainage geometry, land-cover classification, flow measurement, water-quality sampling, parameter estimation, model structure and operational changes. Parameter equifinality means that different parameter combinations can produce similar model outputs. Sensitivity analysis should identify which parameters most affect the decision-relevant outputs, while uncertainty analysis should communicate the range of plausible outcomes rather than presenting a single deterministic result as certain.

Bhusal and Muthukumaran \citeyearpar{bhusal2026} applied SWMM 5.2 to a large, spatially heterogeneous catchment using an EMC-based nutrient approach. The paper is particularly relevant because it reports separate hydrologic and water-quality evaluation. Its hydrology was calibrated and validated, whereas TN evaluation relied on 22 observed concentration values and TP could not be validated owing to missing observations. The authors' EMC sensitivity scenarios improved TN bias but did not improve the reported correlation statistic. This case illustrates that improving average bias is not equivalent to reproducing event-specific concentration dynamics, and that model validation must be pollutant-specific.

\section{GIS representation and catchment discretisation}
GIS supports catchment boundary definition, subcatchment delineation, slope estimation, land-cover mapping, imperviousness assessment and drainage-network construction. DEM resolution, conditioning, coordinate reference system, outlet assignment and ground-truth checks influence the resulting drainage paths and catchment areas. Roofs, roads, parking, landscaped areas and open soil should be separated where their hydrologic or pollutant characteristics justify separate representation.

The appropriate level of discretisation depends on data availability and the decision being supported. Fine spatial resolution can represent source-area heterogeneity but may increase parameter count and calibration burden. Coarse aggregation can hide local source contributions and drainage connectivity. The review will compare studies by catchment scale, number and basis of subcatchment divisions, land-use classes, drainage-network detail and availability of field verification.

\section{Green infrastructure and low-impact development}
GI/LID practices can reduce or delay runoff through infiltration, storage, evapotranspiration, filtration and retention. Candidate measures include green roofs, permeable pavements, bioretention cells, rain gardens, vegetated swales and detention or retention systems. Performance depends on design geometry, soil properties, storage depth, underdrains, antecedent moisture, rainfall regime, clogging and maintenance. EPA-SWMM includes generic LID controls, but the existence of a control object in the software does not establish that its parameterisation or predicted pollutant removal is valid for a particular site \citep{rossman2022}.

Quantity and quality outcomes should be reported separately. A reduction in runoff volume does not necessarily imply a proportional reduction in every pollutant. Treatment performance depends on pollutant form, media properties, bypass, residence time and maintenance. Studies based only on modelled scenarios should be distinguished from field studies with measured inflow and outflow. The completed review will prioritise evidence that reports both hydrologic and water-quality performance and makes clear whether outcomes are observed or simulated.

\section{Critical synthesis and research gaps}
The evidence audited so far supports several provisional conclusions.

\subsection{Integrated evaluation remains less common than hydrologic evaluation}
Niazi et al. \citeyearpar{niazi2017} reported that, among 93 studies with formal evaluation statistics, 67 evaluated runoff and 26 evaluated water-quality endpoints; 22 evaluated both. This imbalance supports a central methodological concern: studies should not be described as fully integrated or fully validated if only runoff has been evaluated.

\subsection{Water-quality observations constrain predictive claims}
Continuous monitoring can reveal event dynamics, but the range of direct high-frequency measurements is limited and proxy relationships require local calibration \citep{razguliaev2024}. Where event observations are sparse, EMC approaches may support approximate load estimation, but they should not be presented as proof of event-scale pollutograph skill \citep{bhusal2026}.

\subsection{Buildup/washoff parameters may be effective rather than universally physical}
Constituent-specific washoff behaviour and catchment-scale parameter transferability require explicit scrutiny. The reviewed evidence does not justify applying one pollutant formulation or parameter set indiscriminately across particulate and dissolved constituents, land uses or catchment scales \citep{gaut2019,bonhomme2017,soonthornnonda2008}.

\subsection{Monsoon transferability requires direct evidence}
The monitoring literature reviewed by Razguliaev et al. \citeyearpar{razguliaev2024} under-represents tropical climates and low- and middle-income settings. Indian case studies should therefore be examined for rainfall intensity and temporal resolution, antecedent dry period, land-use composition, drainage connectivity, pollutant suite, sampling design and QA/QC. Transferability to Bhopal should be argued from these features, not assumed from geographic proximity alone.

\subsection{Provisional research gaps}
The following gaps will be tested against the expanded literature corpus:
\begin{enumerate}[leftmargin=1.5em]
  \item insufficient simultaneous evaluation of runoff hydrographs and pollutant loads/pollutographs;
  \item limited event-based water-quality monitoring and uncertainty characterisation;
  \item uncertain transferability of buildup/washoff parameters across pollutant forms and catchment scales;
  \item inconsistent treatment of first-flush definitions and monitoring resolution;
  \item uncertainty introduced by GIS resolution, source-area delineation and drainage connectivity; and
  \item limited field validation of GI/LID water-quality performance under monsoon conditions.
\end{enumerate}

\section{Proposed framework for a campus catchment in Bhopal}
The proposed application is a future case study, not a completed modelling exercise. The sequence below is intended to make the links between site data, model structure and evaluation explicit.

\begin{enumerate}[leftmargin=1.5em]
  \item \textbf{Verify the study boundary and receiving points.} Confirm the institutional boundary, drainage outfalls, receiving-water connections and any off-site inflows.
  \item \textbf{Build a verified spatial database.} Assemble a suitable DEM, recent land-cover mapping, soil/infiltration information, drainage drawings and field-surveyed network geometry.
  \item \textbf{Characterise rainfall.} Obtain quality-controlled, preferably sub-hourly rainfall records for representative events and document gauge location, missing data and temporal aggregation.
  \item \textbf{Delineate subcatchments and network.} Derive areas, slopes, widths, imperviousness and flow paths using GIS, then verify drainage connectivity in the field.
  \item \textbf{Construct the SWMM quantity model.} Enter verified nodes, conduits, storage elements and outfalls; select routing and infiltration methods with documented rationale.
  \item \textbf{Calibrate and validate hydrology.} Use observed runoff or water-level/discharge data from multiple events, reserving independent events for validation.
  \item \textbf{Select water-quality constituents.} Base selection on a local source inventory, receiving-water concerns and laboratory capacity rather than an arbitrary long constituent list.
  \item \textbf{Design event monitoring.} Pair rainfall and flow with event-based water-quality samples; document sample timing, laboratory methods, QA/QC and detection limits.
  \item \textbf{Evaluate quality representation.} Compare constituent-specific buildup/washoff and, where justified, EMC approaches. Calibrate and validate each pollutant separately where sufficient observations exist.
  \item \textbf{Assess baseline and management scenarios.} Report runoff volume, peak flow, event loads, EMC, pollutographs and first-flush diagnostics, then compare feasible GI/LID scenarios.
  \item \textbf{Quantify uncertainty.} Test sensitivity to rainfall, catchment discretisation, impervious connectivity, pollutant parameters and monitoring error.
\end{enumerate}

The previously discussed approximate campus area and illustrative subcatchment dimensions are not adopted as measured facts in this manuscript. They must be checked against authoritative spatial data and field observations before use. No numerical campus results are reported at this stage.

\section{Limitations and conclusion}
This manuscript is an evolving critical review. The literature search, eligibility screening, full-text retrieval and appraisal are incomplete. The current evidence base includes different levels of access, and abstract-level or publisher-excerpt findings are explicitly distinguished from full-text appraisal. No completed PRISMA flow, pooled performance estimate or formal meta-analysis is claimed. The proposed Bhopal framework remains a research plan pending verified spatial, rainfall, flow and water-quality data.

The literature examined so far indicates that EPA-SWMM offers a practical integrated platform for urban runoff quantity, drainage routing, pollutant buildup/washoff and GI/LID scenario modelling. Its usefulness should not be confused with automatic predictive validity. Integrated credibility requires appropriate catchment and network representation, constituent-specific water-quality assumptions, paired quantity--quality observations, independent validation and transparent uncertainty analysis. EMC methods can be useful when data are limited, but their ability to reproduce event-scale pollutographs must be tested. Buildup/washoff parameters should be interpreted cautiously across pollutant types and spatial scales. A future Bhopal application should therefore be designed around the monitoring and validation requirements from the outset.

\section*{Next manuscript update}
The next update will expand the evidence base in four areas: (i) integrated SWMM calibration and validation, (ii) buildup/washoff parameter uncertainty, (iii) GIS and source-area representation, and (iv) measured GI/LID quantity--quality performance. Sections will be revised only when the new evidence has been extracted and appraised.

\begin{thebibliography}{99}

\bibitem[Bhusal and Muthukumaran(2026)]{bhusal2026}
Bhusal, M. R., and Muthukumaran, S. (2026).
Event Mean Concentration-based Modelling Approach for a Large Scale Spatially Heterogenous Catchment using SWMM.
\textit{Water Resources Management}, 40, Article 55.
\url{https://doi.org/10.1007/s11269-025-04407-z}.

\bibitem[Bonhomme and Petrucci(2017)]{bonhomme2017}
Bonhomme, C., and Petrucci, G. (2017).
Should we trust build-up/wash-off water quality models at the scale of urban catchments?
\textit{Water Research}, 108, 422--431.
\url{https://doi.org/10.1016/j.watres.2016.11.027}.
\textit{Evidence status in this review: abstract/publisher record assessed; full-text appraisal pending.}

\bibitem[Gaut et al.(2019)]{gaut2019}
Gaut, J., Chua, L. H. C., Irvine, K. N., and Le, S. (2019).
Modelling the washoff of pollutants in various forms from an urban catchment.
\textit{Journal of Environmental Management}, 246, 374--383.
\url{https://doi.org/10.1016/j.jenvman.2019.05.118}.
\textit{Evidence status in this review: abstract assessed; full-text appraisal pending.}

\bibitem[Niazi et al.(2017)]{niazi2017}
Niazi, M., Nietch, C., Maghrebi, M., Jackson, N., Bennett, B. R., Tryby, M., and Massoudieh, A. (2017).
Storm Water Management Model: Performance Review and Gap Analysis.
\textit{Journal of Sustainable Water in the Built Environment}, 3(2).
\url{https://doi.org/10.1061/JSWBAY.0000817}.

\bibitem[Page et al.(2021)]{page2021}
Page, M. J., McKenzie, J. E., Bossuyt, P. M., Boutron, I., Hoffmann, T. C., Mulrow, C. D., et al. (2021).
The PRISMA 2020 statement: An updated guideline for reporting systematic reviews.
\textit{BMJ}, 372, n71.
\url{https://doi.org/10.1136/bmj.n71}.

\bibitem[Razguliaev et al.(2024)]{razguliaev2024}
Razguliaev, N., Flanagan, K., Muthanna, T., and Viklander, M. (2024).
Urban stormwater quality: A review of methods for continuous field monitoring.
\textit{Water Research}, 249, 120929.
\url{https://doi.org/10.1016/j.watres.2023.120929}.

\bibitem[Rossman and Simon(2022)]{rossman2022}
Rossman, L. A., and Simon, M. (2022).
\textit{Storm Water Management Model User's Manual, Version 5.2}.
U.S. Environmental Protection Agency, EPA/600/R-22/030.
\url{https://www.epa.gov/system/files/documents/2022-04/swmm-users-manual-version-5.2.pdf}.

\bibitem[Soonthornnonda et al.(2008)]{soonthornnonda2008}
Soonthornnonda, P., Christensen, E. R., Liu, Y., and Li, J. (2008).
A washoff model for stormwater pollutants.
\textit{Science of the Total Environment}, 402(2--3), 248--256.
\url{https://doi.org/10.1016/j.scitotenv.2008.04.036}.
\textit{Evidence status in this review: abstract/publisher excerpts assessed; full-text appraisal pending.}

\end{thebibliography}

\end{document}